% !TeX program = lualatex
\documentclass[a4paper,11pt]{ltjsarticle}
\usepackage{amsmath,amssymb,amsthm}
\usepackage{lmodern}
\usepackage{mathtools}
\usepackage{booktabs}
\usepackage{array}
\usepackage[unicode]{hyperref}

\theoremstyle{definition}
\newtheorem{definition}{定義}[section]
\newtheorem{assumption}[definition]{仮定}
\newtheorem{axiom}[definition]{公理}
\newtheorem{example}[definition]{例}
\newtheorem{remark}[definition]{注記}
\theoremstyle{plain}
\newtheorem{theorem}[definition]{定理}
\newtheorem{proposition}[definition]{命題}

\newcommand{\PoV}{\operatorname{PoV}}
\newcommand{\dfix}{\bigcirc}
\newcommand{\Rel}{\mathbf{Rel}}
\newcommand{\Amphi}{\mathsf{Amphi}}
\newcommand{\QRB}{\mathrm{QRB}}
\newcommand{\Forc}{\mathrm{Forc}}
\newcommand{\sem}[1]{\mathopen{[\![}#1\mathclose{]\!]}}

\title{計算関係論理\\
{\large ---PoV--\texorpdfstring{$\rho$}{rho} インターフェイスによる動的評価体系---}}
\author{千葉将臣}
\date{2026-07-21}

\begin{document}
\maketitle

\begin{abstract}
本稿は、視点数 $\PoV$ と継続残余 $\rho$ を備えた動的関係代数である冗（Amphigory）を全面的な基礎として、計算関係論理（Computational Relational Logic; CRL）を提案する。冗における評価空間、射、交わり、逆関係、残余モジュラー定理、無記、最大不動点 $\nu F$ および外部指示記号 $\dfix$ は既定の構造として継承し、本稿では再定義しない。CRL が新たに導入する中心構造は PoV--$\rho$ インターフェイスであり、冗の抽象的な評価推移を具体的な計算状態間の情報移送として実現する。文脈の融合に伴う視点数の変換と残余の蓄積・圧縮を記述し、冗の双方向演繹が計算モデル上で保存される条件を与える。エルブラン空間上の証明探索を主要例として、順方向、逆方向、極限という三層の計算論的時間を定式化する。最後に、確率的推論および大規模言語モデルとの対応を、同一視ではなく検証を要する解釈仮説として提示する。
\end{abstract}

\section{導入}

Freyd と Scedrov の寓は二項関係の圏 $\Rel$ を抽象化する静的骨格を与える \cite{FreydScedrov1990}。冗はこの寓的骨格に、評価位相、視点数 $\PoV$、継続残余 $\rho$、動的併置、逆方向推論、および評価不能過程を無記として保存する構造を加えた動的関係代数である \cite{ChibaAmphigory2026}。

計算関係論理は、冗の射と評価推移を、証明探索や制約伝播などの具体的な計算状態へ解釈し、異なる評価文脈間で $\PoV$ と $\rho$ を受け渡す規約を PoV--$\rho$ インターフェイスとして付加する。すなわち、冗が動的関係の代数を与え、CRL はその計算的実現と論理的利用法を与える。

本稿の目的は次の四点である。
\begin{enumerate}
 \item 冗の射を具体的な計算状態間の推移として実現する CRL 構造を定義する。
 \item 評価文脈間の情報移送を担う PoV--$\rho$ インターフェイスを定式化する。
 \item 冗の双方向演繹と残余モジュラー定理が計算的実現で保存される条件を示す。
 \item エルブラン空間上の証明探索と確率的推論への応用可能性を示す。
\end{enumerate}

\begin{remark}[本稿の位置づけ]
理論の継承関係は
\[
 \text{Allegory}\longrightarrow\text{Amphigory}
 \longrightarrow\text{CRL}
\]
である。第一の矢印は動的評価構造の付加、第二の矢印は計算的実現とインターフェイス構造の付加を表す。四重残余束（QRB）は、同じ到達不能性の問題を安定テンソル三角圏と強制法の側から幾何学化する \cite{ChibaQRB2026}。CRL は冗から継承した $\dfix$ を用い、QRB との比較も冗を介して行う。
\end{remark}

\section{冗的基盤と PoV--\texorpdfstring{$\rho$}{rho} インターフェイス}

\subsection{冗から継承する構造}

\begin{assumption}[冗的基盤]
$\mathcal{A}$ を \cite{ChibaAmphigory2026} で定義された冗とする。したがって、評価空間
\[
 \mathcal{C}\times\mathcal{V}\times\mathcal{L},
 \qquad
 \mathcal{V}=\mathbb{N}\cup\{\infty\},
 \qquad
 \mathcal{L}=\mathbb{R}_{\geq0},
\]
評価付きの射 $R:X\to_{\langle p,\rho\rangle}Y$、交わり $R\cap S$、逆関係 $R^\circ$、合成、残余モジュラー定理、および無記に関する極限構造が与えられているものとする。本稿におけるこれらの記号と法則は、すべて冗における定義と同一であり、CRL の射とは、計算的解釈を与えられた冗の射にほかならない。
\end{assumption}

\subsection{CRL 構造}

\begin{definition}[冗上の CRL 構造]
冗 $\mathcal{A}$ 上の CRL 構造とは、次のデータからなる。
\begin{enumerate}
 \item 各対象 $X$ に対する計算状態の解釈 $\mathsf{State}(X)$。
 \item 各冗の射 $R:X\to_{\langle p,\rho\rangle}Y$ に対する状態推移関係 $\sem{R}\subseteq\mathsf{State}(X)\times\mathsf{State}(Y)$。
 \item 合成可能または併置可能な射の評価情報を接続する PoV--$\rho$ インターフェイス $\Phi$。
 \item 各状態推移に対する実現残余 $\widehat{\rho}(\sem{R})$。これは冗の抽象残余 $\rho(R)$ の計算モデル上の測定値である。
\end{enumerate}
状態推移の解釈は、冗の合成、交わり、逆関係、および評価推移を保存するものとする。
\end{definition}

\subsection{インターフェイス}

\begin{definition}[PoV--$\rho$ インターフェイス]
冗の射 $R,S$ の評価情報を接続する PoV--$\rho$ インターフェイスとは、部分写像
\[
 \Phi_{R,S}:
 (\mathcal{V}\times\mathcal{L})^2
 \dashrightarrow
 \mathcal{V}\times\mathcal{L}
\]
であって、二つの評価文脈から接続後の視点数と残余を生成するものをいう。冗の動的併置と整合する基本形を
\[
 \Phi_{R,S}\bigl((n,\alpha),(m,\beta)\bigr)
 =\bigl(n+m,\alpha+\beta+\Delta_{R,S}(n,m)\bigr)
\]
とする。$\Delta_{R,S}(n,m)\geq0$ は文脈干渉項である。
\end{definition}

インターフェイスは冗の交わりを再定義するものではなく、冗に既に存在する干渉項を具体的な計算モデルでどのように測定し、次の状態へ渡すかを定める規約である。$\Delta=0$ は独立な文脈の併置を、$\Delta>0$ は共有変数、競合制約、観測順序などによる干渉を表す。

\begin{assumption}[インターフェイスの整合性]
$\Phi$ は冗の恒等射、合成および交わりと整合し、三つの文脈を順に接続する二通りの評価が一致するものとする。すなわち、定義される範囲で単位律と結合律を満たし、冗で計算された $\PoV$ と $\rho$ を保存する。
\end{assumption}

\section{冗の双方向演繹の計算的実現}

\subsection{継承される演算}

冗において、交わりは観測文脈の動的併置であり、逆関係は未来制約を現在の文脈へ引き戻す逆方向推論である。CRL はこれらを新たな演算として定義せず、次の既存の評価法則をそのまま用いる。
\[
 \PoV(R\cap S)=\PoV(R)+\PoV(S),
\]
\[
 \rho(R\cap S)
 =\rho(R)+\rho(S)+\Delta_{R,S}
\]
\[
 \PoV(R^\circ)=\PoV(R)
 .
\]

CRL に固有なのは、これらの抽象演算を状態推移 $\sem{R}$ とインターフェイス $\Phi$ によって実行可能な形へ移すことである。逆関係そのものを時間反転とはみなさず、後段で判明した制約を先行状態へ伝播して未展開の分岐を削除する操作的順序として実現する。

\subsection{残余モジュラー定理}

冗の残余モジュラー定理は
\[
 RS\cap T\subseteq R(S\cap R^\circ T),
 \qquad
 \rho(RS\cap T)\geq\rho\bigl(R(S\cap R^\circ T)\bigr)
\]
を与える。CRL モデルでは、この抽象的な残余比較が実現残余 $\widehat{\rho}$ による具体的な探索費用の比較として保存されることを要求する。

\begin{assumption}[局所制約による残余圧縮]
任意の適合する射 $R,S,T$ に対して、未来制約 $T$ を $R^\circ T$ として先に引き戻す評価戦略は、事後的に $T$ と照合する戦略より残余を増加させないとする。すなわち
\[
 \widehat{\rho}\bigl(\sem{R(S\cap R^\circ T)}\bigr)
 \leq
 \widehat{\rho}\bigl(\sem{RS\cap T}\bigr)
\]
が成り立つ。
\end{assumption}

\begin{theorem}[残余モジュラー定理の CRL 保存]
冗上の CRL 構造がインターフェイスの整合性と局所制約による残余圧縮を満たすならば、その状態推移の解釈において
\[
 \sem{RS\cap T}\subseteq\sem{R(S\cap R^\circ T)},
\]
\[
 \rho_{\mathrm{left}}\geq\rho_{\mathrm{right}}
\]
が成り立つ。ここで
\[
 \rho_{\mathrm{left}}:=\widehat{\rho}\bigl(\sem{RS\cap T}\bigr),
 \qquad
 \rho_{\mathrm{right}}:=\widehat{\rho}\bigl(\sem{R(S\cap R^\circ T)}\bigr)
\]
である。
\end{theorem}

\begin{proof}
抽象的な包含と残余不等式は冗の残余モジュラー定理から継承される。状態推移の解釈は冗の演算を保存するため、第一式が得られる。PoV--$\rho$ インターフェイスは評価情報を保存し、局所制約による実現残余の圧縮を仮定しているため、第二式も具体的な CRL モデル上で成立する。
\end{proof}

\begin{remark}
抽象的な残余不等式は冗の定理であるが、特定の実装で測定した実行時間やメモリ量が同じ不等式を満たすとは限らない。制約の引き戻し自体が高価な場合には、その実装は上の圧縮仮定を満たす CRL モデルではない。
\end{remark}

\section{冗から継承される無記と外部指示}

CRL は、非停止評価に関する新たな公理や外部点を導入しない。冗の無記公理により、評価列に沿って
\[
 \PoV\longrightarrow\infty
\]
となる過程は任意の通常値へ同一視されず、無記として残余空間へ隔離される。完備束 $L$ 上の単調作用素 $F$ に対する内部表現は、冗で定義された中道不動点
\[
 \nu F:=\bigvee\{x\in L\mid x\leq F(x)\}
\]
であり、$\dfix$ はその値や $L$ の追加要素ではない。指示関係
\[
 \dfix\rightsquigarrow\nu F
\]
も冗からそのまま継承する。

CRL の役割は、この極限を証明探索の循環、無限分岐、資源制約下の未決定などの具体的な計算過程として解釈することにある。したがって、CRL の計算モデルが $\dfix$ を生成するとは、冗の外部指示を計算状態上で観測したことを意味し、新しい値を計算したことを意味しない。

\begin{remark}[QRB との対応]
CRL が用いる $\dfix$ は冗の指示記号そのものである。したがって QRB における強制法の指示記号 $\dfix_{\Forc}$ および QRB の指示記号 $\dfix_{\QRB}$ との比較も、冗で述べられた対応を継承する \cite{ChibaAmphigory2026,ChibaQRB2026}。CRL は $\nu F$、ジェネリック・フィルター、QRB の代表点の新たな同一視を主張しない。
\end{remark}

\section{エルブラン空間上の計算論的時間}

\subsection{エルブラン・関係}

一階シグネチャ $\Sigma$ のエルブラン宇宙を $H_\Sigma$ とする。論理プログラムのゴール、代入、制約集合を状態として含む集合を $\mathsf{State}(H_\Sigma)$ と書く。SLD 導出および制約論理プログラミングの標準的背景については \cite{Lloyd1987,JaffarLassez1987} を参照する。

\begin{definition}[エルブラン・関係]
証明探索の一段階を表す CRL 射
\[
 R:\mathsf{State}(H_\Sigma)
 \longrightarrow_{\langle p,\rho\rangle}
 \mathsf{State}(H_\Sigma)
\]
をエルブラン・関係と呼ぶ。$p$ は選択可能な節、リテラル、代入候補による分岐量を、$\rho$ は失敗した単一化、未解決制約、再訪した状態の重み付き総和を表す。
\end{definition}

\begin{proposition}[計算論的時間の三層]
エルブラン・関係による探索は、次の三種類の操作的順序を持つ。
\begin{enumerate}
 \item 順方向時間：ゴール簡約と部分ゴール生成による探索状態の更新。
 \item 逆方向時間：失敗情報、型制約、終端条件のバックトラックまたは制約伝播。
 \item 極限時間：循環または無限分岐により有限の $\PoV$ では決定できず、無記として隔離される評価列。
\end{enumerate}
\end{proposition}

\begin{proof}
第一項は解像ステップの向き、第二項は失敗した後続状態から先行状態への制約移送に対応する。第三項は冗から継承した無記公理を、循環または無限分岐を持つ評価列へ適用したものである。
\end{proof}

\begin{example}[循環を含む経路探索]
有向グラフ上の述語 $\mathsf{path}(x,y)$ を再帰的な Horn 節で定義する。循環を含むグラフで $\mathsf{path}(a,z)$ を深さ優先探索すると、節と辺の選択が $\PoV$ を増加させ、失敗した単一化と再訪状態が $\rho$ に蓄積される。到達先 $z$ から逆向きに到達可能性制約を伝播すれば分岐を削減できるが、その圧縮量は探索戦略とグラフ構造に依存する。無限循環が検出も切断もされない場合、その評価列は通常の成功・失敗へ還元せず、$\dfix$ によって外部化を指示する。
\end{example}

\section{確率的推論への解釈仮説}

CRL の評価位相 $\mathcal{C}$ は、多値または確率的な評価を収容できる。ただし、$\PoV$ や $\rho$ を既存の機械学習パラメータと直接同一視することはできない。本節では、検証対象となる解釈仮説だけを述べる。確率的推論の一般的背景については \cite{Bishop2006}、アテンション機構については \cite{Vaswani2017} を参照する。

\begin{assumption}[確率的解釈仮説]
確率的推論過程に対して、次の量を観測可能にできると仮定する。
\begin{enumerate}
 \item 有効視点数 $\PoV_{\mathrm{eff}}$：同時に保持される相互に区別可能な仮説・文脈・サンプルの有効数。
 \item 確率的残余 $\rho_{\mathrm{prob}}$：棄却された仮説、未説明分散、制約違反、計算予算超過を重み付けした量。
 \item 干渉項 $\Delta$：文脈融合による情報重複または矛盾の追加費用。
\end{enumerate}
\end{assumption}

\begin{center}
\begin{tabular}{>{\raggedright\arraybackslash}p{0.30\textwidth} >{\raggedright\arraybackslash}p{0.56\textwidth}}
\toprule
CRL の量 & 確率的推論における候補的解釈 \\
\midrule
$\PoV$ & 有効サンプル数、競合仮説数、保持された文脈の多様性 \\
$\rho$ & 未説明損失、棄却分岐、制約違反、計算予算の残差 \\
$R^\circ$ & 後段の損失・制約を前段表現へ伝播する逆方向処理 \\
$\dfix$ & 所定の計算予算と観測文脈では決定できないことの指示 \\
\bottomrule
\end{tabular}
\end{center}

大規模言語モデルにおいて、アテンション・ヘッド数、温度、コンテキスト長をそのまま $\PoV$ と置くことはしない。それらから推定される有効な仮説多様性を $\PoV_{\mathrm{eff}}$ の候補とする。同様に、$\rho$ を幻覚の度合いそのものとはせず、事実制約違反、校正誤差、棄却候補、計算費用を含む複合指標として設計する必要がある。

この仮説を理論へ昇格させるには、具体的な推論アルゴリズムごとに状態推移関係 $\sem{R}$、$\Phi_{R,S}$、$\Delta$ を定義し、残余圧縮仮定が成立する条件を実験的・数学的に検証しなければならない。

\section{結論}

本稿は、冗を全面的な基礎とし、その射と評価推移を具体的な計算状態へ解釈する計算関係論理を提案した。評価空間、PoV 付きの射、交わり、逆関係、残余モジュラー定理、無記、最大不動点および外部指示は冗から継承し、CRL 内では再定義しなかった。CRL 固有の構造は、状態推移の解釈と PoV--$\rho$ インターフェイスである。

CRL の中心的な特徴は次の三点である。
\begin{enumerate}
 \item 冗の抽象的な射を、証明探索や制約伝播の状態推移関係として実現する。
 \item 冗の $\PoV$、$\rho$ および干渉項を、PoV--$\rho$ インターフェイスによって計算文脈間で受け渡す。
 \item 冗から継承した双方向演繹と外部指示を、エルブラン探索および確率的推論の共通語彙として用いる。
\end{enumerate}

今後の課題は、具体的モデルにおける状態推移関係、$\Delta$ と $\rho$ の測定法、インターフェイスの単位律・結合律、残余圧縮仮定の十分条件を与えることである。QRB との比較は、冗で構成された対応を基礎として計算モデルへ持ち上げる必要がある。これらが確立されれば、CRL は冗の動的関係代数を具体的な計算意味論へ接続する論理体系として位置づけられる。

\bibliographystyle{plain}
\bibliography{../ref}

\end{document}
